In parametric Bayesian statistics, one works with statistical models, namely probability distributions indexed by finite-dimensional parameters. In many applications, this assumption can be restrictive. Bayesian nonparametric statistics therefore provides more flexible modeling approaches. The fundamental object to work with is a random probability measure, namely a random variable taking values in the space of probability measures that may have generated the observed data. Before formally defining this object, we introduce the setting for the statistical models.

Let \(\X\) be a Polish space equipped with its Borel \(\sigma\)-algebra \(\mathcal{B}(\X)\). Denote by $\probx$ the space of probability measures on \((\X,\mathcal{B}(\X))\). We equip $\probx$ with the Borel $\sigma$-algebra generated by the topology of weak convergence. This $\sigma$-algebra is the smallest one that makes the maps $\phi_A : \probx \to [0,1]$, defined by
\[
\phi_A(P) = P(A),
\]
measurable for all \(A \in \mathcal{B}(\X)\) (see Proposition A.5 in \cite{ghosal2017fundamentals}).

\begin{definition}\label{definition:random_probability_measure}
Let $\probspace$ be a probability space. Random probability measure is a measurable map from $\Omega$ to $\probx$.
\end{definition}

In the literature, random probability measures are also often defined in terms of transition probability kernels. Under the measurable structure considered here, these two formulations are equivalent. We will denote the random probability measure by $\tilde{P}.$

A fundamental idea in Bayesian statistics is that, to make a statistical inference, the order in which the observations are observed should not matter. This symmetry is formalized by the notion of exchangeability, introduced by de Finetti.
\begin{definition}\label{definition:exchangeability}
A sequence of random variables $(X_i)_{i \geq 1}$ is exchangeable if
\[
\big(X_1,\dots, X_n\big) \overset{d}{=} \big(X_{\sigma(1)},\dots, X_{\sigma(n)}\big),
\]
for all $n \geq 1$ and every permutation $\sigma$ of $\{1,\dots,n\}$.
\end{definition}
By de Finetti's representation theorem, exchangeability is equivalent to the assumption that the law of the observations is the mixture of product probability measures weighted by the prior distribution. 
\begin{theorem}\label{theorem:definetti}
A sequence of random variables $(X_i)_{i \geq 1}$ is exchangeable if and only if there exists a probability measure $Q \in \probprobx$ such that
\[
\mathbb{P}\left(X_1\in A_1,\dots, X_n\in A_n\right) = \int_{\probx}\prod_{i = 1}^n P(A_i)\, \mathrm{d}Q(P)
\]
for all $n \geq 1$ and $A_1,\dots,A_n \in \mathcal{B}(\X)$.
\end{theorem}
A consequence of the above theorem is the following interpretation of the data generation process:
\begin{align*}
    &\tilde{P} \sim Q \\
    &X_1,\dots X_n|\tilde{P}\overset{\iid}{\sim} \tilde{P},
\end{align*}
where $Q \in \probprobx$ is referred as prior distribution. 

In Bayesian nonparametrics, one of the early central challenges was to construct prior distributions on infinite-dimensional parameter spaces whose posterior distributions remained analytically or computationally tractable. A major breakthrough in this direction was the introduction of the Dirichlet process in \cite{ferguson1973bayesian}. The Dirichlet process provides a prior distribution over random probability measures and possesses a conjugacy property analogous to that of finite-dimensional Bayesian models.

The method used by Ferguson to construct the Dirichlet process relies on finite-dimensional distributions. Specifically, one specifies the joint distribution of the probabilities assigned to finite measurable partitions and verifies the consistency conditions required to guarantee the existence of a random probability measure. This approach is analogous to the classical construction of stochastic processes through their finite-dimensional distributions. We now briefly describe this method.

Recall that, we aim to define a random variable $\tilde{P}$ taking values in the space of functions from $\mathcal{B}(\X)$ to $[0,1]$, endowed with the additional properties required of a probability measure. The main tool for this construction is the Kolmogorov extension theorem. This result allows one to define a stochastic process by specifying a collection of finite-dimensional distributions that satisfy suitable consistency conditions. In our setting, these finite-dimensional distributions are indexed by finite collections of measurable sets $A_1, \dots, A_n \in \mathcal{B}(\X)$. It is sufficient, however, to define these laws on measurable partitions of the space (see page 213 of \cite{ferguson1973bayesian}).

\begin{definition}\label{definition:measurable_partition}
A measurable partition of a measurable space $(\X, \mathcal{B}(\X))$ is a collection of sets $B_1, \dots, B_k \in \mathcal{B}(\X)$ such that the sets are pairwise disjoint and
\[
\X = \bigcup_{i=1}^k B_i.
\]
\end{definition}

The first requirement is that the random measure assigns zero mass to the empty set, namely
\[
\law\bigl(\tilde{P}(\emptyset)\bigr) = \delta_0.
\]

The second requirement is the consistency condition. This condition states that if a measurable partition is refined into a finer partition, then the induced joint distributions must agree. More precisely, consider two measurable partitions $(B_1, \dots, B_k)$ and $(C_1, \dots, C_m)$, where each set $B_i$ can be expressed as a union of consecutive elements of the finer partition: For all $i = 1,\dots, k,$
\[
B_i = \bigcup_{j=r_i}^{r_{i+1}-1} C_j,
\]
with indices satisfying $1 = r_1 < r_2 < \dots < r_{k+1} = m+1$. Then the consistency condition requires that
\[
\law\left(
\sum_{j=r_1}^{r_2-1} \tilde{P}(C_j), \dots,
\sum_{j=r_k}^{r_{k+1}-1} \tilde{P}(C_j)
\right)
=
\law\bigl(\tilde{P}(B_1), \dots, \tilde{P}(B_k)\bigr).
\]

Under these conditions, Lemma 1 of \cite{ferguson1973bayesian} guarantees the existence of a random probability measure whose finite-dimensional distributions coincide with those specified in the construction.